\documentclass[11pt]{article}
\usepackage{amsmath, amssymb}
\usepackage[margin=1in]{geometry}
\usepackage{booktabs}
\usepackage{graphicx}
\PassOptionsToPackage{obeyspaces,spaces}{url}
\usepackage{url}
\usepackage[hidelinks]{hyperref}
\setlength{\parskip}{4pt}
\setlength{\parindent}{0pt}
\graphicspath{{model_output/moist_v1_tauc_dz/}}

\newcommand{\py}{\partial_y}
\newcommand{\Shat}{\hat S}
\newcommand{\Hhat}{\hat H}
\newcommand{\Wq}{W_q}
\newcommand{\Wstar}{W^*}
\newcommand{\Hv}{\mathcal{H}}

\title{Moist V1 status report: the parameter-sensitivity program\\
and the gross-moist-stability regime boundary}
\date{2026-07-19}

\begin{document}
\maketitle

\section*{Summary}

Moist V1 (a passive column-water-vapor field $W$ riding on the dry
Sobel--Schneider circulation) is implemented, validated, and its full
placeholder-parameter sensitivity program is complete: four scans covering
the eddy moisture diffusivity $D$, the precipitation threshold $W_c$, the
convective timescale $\tau_c$, and the stratification depth $\Delta_z$,
eighteen steady-state runs in all. One result organizes everything
(Fig.~\ref{fig:regime}): because the equilibrium overturning is weak, $W$ is
nearly uniform at the value $\Wq = W_c + \tau_c E_0$ set by the local
evaporation--precipitation balance, so the sign of the gross moist stability
$\Hhat = \Shat - L_v(2a{-}1)W$, and with it the direction of the mean
meridional MSE flux, is decided by a single comparison: $\Wq$ against
$\Wstar = \Shat/(L_v(2a{-}1))$. Every crossover this comparison predicts in
closed form is confirmed by the runs. At the
boundary the mean MSE flux vanishes by exact cancellation of its dry-static
and latent parts, leaving the diagnostic eddy flux $-L_v D\,\py W$ as the
only meridional MSE transport: the regime the eddy-diffusion program is
aimed at, now locatable exactly. Remaining before V2: the ERA5 calibration
of $a$ and $D$ (Spencer), then latent heating.

\section{Where the moist model stands}

V1 adds one prognostic equation to the dry model: column water vapor
$W(y,t)$, transported by the lower-branch flow and mixed down-gradient,
with uniform evaporation $E_0$ and a quasi-equilibrium precipitation sink
$P = (W - W_c)^+/\tau_c$ (Appendix~\ref{app:eqs} has the full equation
set). The coupling is strictly one-way: $W$ feeds back on nothing, and the
dry fields of a moist run are bit-for-bit identical to a dry twin, a tested
invariant that also holds across every run discussed here. V2 will close
the loop by letting $P$ heat the temperature equation.

Implementation status: merged to \texttt{main}, 338 tests, with exact
invariants (moisture mass conservation to roundoff, mirror parity, bit-exact
restarts, dry invariance) and a regression baseline. The moist parameters
are placeholders pending an ERA5 calibration of the lower-layer moisture
fraction $a$ and the diffusivity $D$; that calibration is the only open
pre-V2 task besides V2 itself.

The sensitivity program ran on the production dry formulation (staggered
grid, EMFD gate on, MC-limited stencil; \texttt{ny}=801, $\Delta t=30$~s,
equinoctial $\sin^2$ forcing) with every run integrated to a slow-drift-gated
statistical steady state (gate firings: day 1772--1896). The four scans:
\begin{center}
\begin{tabular}{llll}
\toprule
scan & values & runs & center point \\
\midrule
$D$ & $\{0, 1, 2\}\times10^6$ m$^2$/s & 3 & shared: the default run \\
$W_c$ & 35, 40, 44, 50, 55 kg/m$^2$ & 4 (+center) & ($W_c$=50) \\
$\tau_c$ & 0.5, 2, 4, 16, 48 h; pair at $W_c$=40 & 6 (+center) & ($\tau_c$=4 h) \\
$\Delta_z$ & 30, 45, 60, 68.2, 75, 90 K & 5 (+center) & ($\Delta_z$=60 K) \\
\bottomrule
\end{tabular}
\end{center}
All other parameters sit at their defaults (Appendix~\ref{app:params});
the default run (\mbox{$W_c$=50}, $\tau_c$=4~h, $\Delta_z$=60~K, $D$=$10^6$)
is the shared center of every ladder.

\section{One boundary organizes all four dials}
\label{sec:boundary}

\begin{figure}[tbp]
\centering
\includegraphics[width=0.8\linewidth]{moist_v1_regime_plane.png}
\caption{Every scanned run placed by what its parameters set: the quiescent
plateau $\Wq = W_c + \tau_c E_0$ (the $W$ the moist closure relaxes to where
the circulation is weak) against the crossover $\Wstar = \Shat/(L_v(2a{-}1))$
(the $W$ at which $\Hhat$ changes sign, proportional to $\Delta_z$). Fill
color gives the measured sign of $\Hhat(y)$ over the whole domain. Runs sit
on the predicted side of the diagonal in every case; the one mixed-sign run
($\Delta_z$=68.2~K) sits on it. The $D$ ladder collapses onto the shared
center point: $D$ moves neither coordinate.}
\label{fig:regime}
\end{figure}

The equilibrium overturning is weak because it is eddy-driven: the zonal
momentum balance is $\beta y\,v \approx$ EMFD (drag is two orders down),
giving $|v|_{\max} = 0.37$~m/s at the default $v_d = 2.5$~m/s. Weak
transport means the moisture equation is nearly in local balance
everywhere: $P \approx E_0$, so $W$ sits within 0.01--1.4~kg/m$^2$ of the
\emph{quiescent plateau}
\begin{equation}
\Wq \;=\; W_c + \tau_c E_0,
\end{equation}
the value at which the precipitation closure removes exactly what
evaporation supplies. (Measured: plateau matches this to $\le 7.5\times
10^{-8}$~kg/m$^2$ in every run.)

With $W$ pinned near $\Wq$, the gross moist stability
$\Hhat(y) = \Shat - L_v(2a{-}1)\,W(y)$ has, to the accuracy of those small
departures, one sign for the whole domain, decided by $\Wq$ against
\begin{equation}
\Wstar \;=\; \frac{\Shat}{L_v(2a{-}1)}
      \;=\; \frac{C\,\delta\,\Delta_z/H}{L_v(2a{-}1)}
      \;=\; 44.57~\text{kg/m}^2 \times \frac{\Delta_z}{60~\text{K}} .
\end{equation}
$\Hhat > 0$ ($\Wq < \Wstar$): the mean overturning exports MSE from the
ITCZ, the familiar dry-side regime. $\Hhat < 0$: the mean MSE flux is
up-gradient (into the ITCZ), the moisture-mode side. Figure~\ref{fig:regime}
shows the runs in the $(\Wq, \Wstar)$ plane; the diagonal is the
regime boundary, and each dial moves the system along exactly one axis:
$W_c$ and $\tau_c$ move $\Wq$, $\Delta_z$ moves $\Wstar$, $D$ moves
neither.

Three crossover predictions follow, and the runs confirm each
(all numbers are $\Hhat$ at the ITCZ in MJ/m$^2$):
\begin{itemize}
\item \textbf{In $W_c$} at fixed $\tau_c$: $W_c^* = \Wstar - \tau_c E_0 =
43.91$. Measured ladder $+15.5, +6.7, -0.3, -10.8, -19.6$ across
$W_c = 35 \ldots 55$; the $W_c$=44 run is in effect the boundary run.
\item \textbf{In $\Delta_z$}: $\Delta_z^* = 68.20$~K at the default $\Wq$.
The run placed at 68.2~K lands mixed-sign with $\Hhat(0) = -0.16$.
\item \textbf{In $\tau_c$}: only reachable if $W_c < \Wstar$, at
$\tau_c^* = (\Wstar - W_c)/E_0$. At $W_c$=40 that is $9.94\times10^4$~s
(1.15~d), and the run pair at $\tau_c$ = 0.5~h / 2~d brackets it:
$\Hhat(0) = +7.84 \to -7.88$. At the default $W_c$=50 the plateau starts
above $\Wstar$, so shrinking $\tau_c$ only saturates $\Hhat(0)$ at
$\Shat - L_v(2a{-}1)W_c = -9.5$; no $\tau_c$ crosses.
\end{itemize}

\textbf{At the boundary the eddy flux is the only MSE transport.} The mean
flux $v\Hhat$ is the residual of two large opposing parts, the poleward
dry-static flux $\Shat v$ and the equatorward latent flux $-L_v(2a{-}1)Wv$.
At $\Delta_z$=68.2 these are $+30.18$ and $-30.18$~MW/m at the flux
latitude: the net is $0.00$, and the diagnostic eddy flux
($0.13$~MW/m) carries everything. The $W_c$ scan found the same at its
crossover (eddy twice the net). Away from the boundary the mean part
dominates by factors of 20--60.

\textbf{The mixed-sign window is as narrow as the $W$ field is flat.} The
spatial spread of $\Hhat$ is $L_v(2a{-}1)$ times the spatial spread of $W$
($\approx 0.17$~kg/m$^2$ at default $\tau_c$), about $0.3$~MJ/m$^2$.
Turning any dial slides the nearly flat $\Hhat(y)$ rigidly through zero, so
both signs coexist over only $\sim$0.17~kg/m$^2$ of $W_c$ or
$\sim$0.23~K of $\Delta_z$: the domain changes regime almost as a unit.
The window widens in proportion to the transport departures, hence linearly
in $\tau_c$ ($\sim$3~K of $\Delta_z$ at $\tau_c$=2~d).

Finally, every scanned configuration precipitates everywhere
($P_{\min} \ge 3.4$~mm/d against $E_0 = 3.97$~mm/d): with uniform
evaporation and departures this small, $W$ never falls below $W_c$
anywhere. A subtropical dry zone would need either structured evaporation
or a much stronger circulation.

\section{The convective timescale $\tau_c$}
\label{sec:tauc}

\begin{figure}[tbp]
\centering
\includegraphics[width=\linewidth]{moist_v1_tauc_scan.png}
\caption{$\tau_c$ ladder at $W_c$=50 (green $\to$ purple: 0.5~h $\to$ 2~d).
Clockwise from top left: $W$ (the plateau climbs as $W_c + \tau_c E_0$);
$P$ (nearly unchanged: it is fixed by the moisture budget); departures
$W - \Wq$ (amplitude grows with $\tau_c$, shape fixed); $\Hhat$ (deepens,
no sign change at this $W_c$).}
\label{fig:tauc-ladder}
\end{figure}

\begin{figure}[tbp]
\centering
\includegraphics[width=\linewidth]{moist_v1_tauc_scan_summary.png}
\caption{$\tau_c$ summary. Top left: the departure shape
$(W-\Wq)/\tau_c = P - E_0 = -\py F$ collapses across the ladder: the
transport convergence is $\tau_c$-invariant while the departure amplitude
scales with $\tau_c$. Top right: that amplitude scaling (log--log, slope-1
guide). Bottom left: $\Hhat(0)$ against $\tau_c$ for $W_c$=50 (saturates)
and $W_c$=40 (crosses zero at the predicted $\tau_c^*$, dashed); points sit
below the plateau-based curves by the ITCZ-bump contribution
$L_v(2a{-}1)\,\Delta W$. Bottom right: mean-flux decomposition at the flux
latitude.}
\label{fig:tauc-summary}
\end{figure}

$\tau_c$ has one trivial role and one dynamical role.

\textbf{Trivial: it sets the plateau.} A column raining at rate $P$ must
hold $P\tau_c$ of water above threshold, so where $P = E_0$ the column
holds $W_c + \tau_c E_0$. Slower convective removal means a moister column
at the same rain rate. This is a single-point balance with no dynamics in
it; the runs track it to $\le 7.5\times10^{-8}$~kg/m$^2$ over a
$\times$96 range of $\tau_c$ (0.5~h to 2~d).

\textbf{Dynamical: it is the gain on the transport structure.} At steady
state $(W - \Wq)/\tau_c = P - E_0 = -\py F$ exactly, where $F$ is the total
meridional moisture flux. The measured departure profiles collapse under
this scaling (Fig.~\ref{fig:tauc-summary}, top left): the moisture-flux
convergence is set by the dry circulation and the nearly fixed $W$, and is
therefore $\tau_c$-invariant, while the $W$ departures needed to convert
that convergence into extra or reduced rain scale with $\tau_c$. Measured:
the ITCZ bump grows $0.013 \to 1.359$~kg/m$^2$ across the ladder, with
$(\max W - \Wq)/\tau_c$ constant to within 11\% (the excess tracks the
plateau's own growth, since $F \propto W$). Equilibrium $P$ itself moves by
at most $0.09$~mm/d across the whole ladder. Consequences: $\Hhat(0)$
deepens linearly, $-9.7 \to -25.8$~MJ/m$^2$ at $W_c$=50, and the mixed-sign
window of Sec.~\ref{sec:boundary} widens linearly.

\textbf{Regime reach.} As shown in Sec.~\ref{sec:boundary}: crossing in
$\tau_c$ requires $W_c < \Wstar$ (confirmed at $W_c$=40); at the default
$W_c$=50 the dial saturates on the moisture-mode side.

\textbf{Numerics.} $D$ and $P$ are stepped at the lagged $n{-}1$ level, so
the precipitation stability bound is $2\Delta t/\tau_c \le 1$, i.e.\
$\tau_c \ge 60$~s at $\Delta t$=30. Ten-day smokes at $\tau_c = 60$, 120,
300~s (bound ratios 1.0, 0.5, 0.2) all ran clean, so the bound is not a
practical constraint for any physical $\tau_c$. The dry fields of every
$\tau_c$ run (including the $W_c$=40 pair) are bit-identical to the center
run ($\max|\Delta u| = 0$ over all 1855 daily fields; all gate day 1855),
as passivity requires.

\section{The stratification depth $\Delta_z$}
\label{sec:dz}

\begin{figure}[tbp]
\centering
\includegraphics[width=\linewidth]{moist_v1_dz_scan.png}
\caption{$\Delta_z$ ladder (green $\to$ purple: 30 $\to$ 90~K). $W$ and $P$
keep the same plateau (the moist closure does not see $\Delta_z$) with
structure amplitude following the weakening dry $v$; $\Hhat$ slides through
zero; the mean MSE flux $v\Hhat$ reverses direction across the ladder.}
\label{fig:dz-ladder}
\end{figure}

\begin{figure}[tbp]
\centering
\includegraphics[width=\linewidth]{moist_v1_dz_scan_summary.png}
\caption{$\Delta_z$ summary. Top left: $\Hhat(0)$ crosses zero at the
predicted $\Delta_z^* = 68.2$~K. Top right: flux decomposition; the DSE
flux rises with $\Shat$ but sublinearly (the circulation weakens), the
latent flux shrinks with $v$, the net reverses at the crossover, the eddy
flux stays at 0.1--0.25 MW/m throughout. Bottom left: dry-circulation response (percent of
the $\Delta_z$=60 values). Bottom right: the regime diagram at $W_c$=50 in
the $(\tau_c, \Delta_z)$ plane with the boundary
$\Delta_z^*(\tau_c)$ and the mixed-sign band from the measured $W$ ranges.}
\label{fig:dz-summary}
\end{figure}

$\Delta_z$ is the one dial with a dual role: it sits in the dry
thermodynamic equation (the adiabatic term $(\delta\Delta_z/H)\py v$), so it
changes the circulation itself, and it sets the gross dry stability
$\Shat = C\delta\Delta_z/H$, so it moves the crossover $\Wstar$
proportionally. The moist closure does not see it: the plateau stays at
$\Wq = 50.66$ across the whole ladder (measured ITCZ $W$: 50.74--50.83 over
a $\times$3 range of $\Delta_z$).

\textbf{Energetics.} $\Hhat(0)$ runs from $-49.9$~MJ/m$^2$ at 30~K to
$+28.2$ at 90~K, crossing zero at the predicted 68.20~K; the mean MSE flux
at the flux latitude runs $-26.5 \to +8.3$~MW/m, reversing there. The DSE
component grows only 63\% ($20.9 \to 34.1$~MW/m) over a $\times$3 growth in
$\Shat$ because the circulation weakens at the same time; the latent
component shrinks with $|v|$ ($-47.4 \to -25.8$~MW/m) at fixed $W$.

\textbf{Dry response.} Over 30~$\to$~90~K the jet weakens $31.9 \to
25.9$~m/s ($-19$\%) and $|v|_{\max}$ falls $0.53 \to 0.29$~m/s ($-45$\%),
with the flux latitude moving poleward ($2.56 \to 3.03$~Mm) and the
convergence gate firing progressively later (day 1772 $\to$ 1896).
Mechanism: the steady thermodynamic balance
$(\delta\Delta_z/H)\py v = (\theta_E - \theta)/\tau$ requires less
divergence per unit diabatic heating when the stratification is deeper,
so $v$ weakens with $\Delta_z$. The measured response is shallower than the
fixed-heating $1/\Delta_z$ scaling ($\times$1.8 over a $\times$3 range)
because the disequilibrium $\theta_E - \theta$ adjusts as well.

\section{The other two dials: $W_c$ and $D$}
\label{sec:wcd}

The $W_c$ and $D$ scans (runs and figures in
\path|model_output/moist_v1_wc_scan/| and
\path|model_output/moist_v1_validation/|):

\textbf{$W_c$ (35--55).} Shifts $\Wq$ one-for-one. The $\Hhat(0)$ ladder
runs from $+15.5$ to $-19.6$~MJ/m$^2$ across $W_c = 35 \ldots 55$, crossing
at $W_c^* = 43.91$; the $W_c$=44 run sits at the boundary
($\Hhat(0) = -0.32$, net mean flux $-0.06$~MW/m, eddy flux $0.12$, i.e.\
twice the net).

\textbf{$D$ (0, 1, 2 $\times 10^6$~m$^2$/s).} In V1, $D$ is the eddy-flux
dial and nothing else: the eddy MSE flux scales with it (0, 0.28,
0.46~MW/m), while its leverage on $W$ itself is negligible
($\max|W_{D2} - W_{D0}| = 0.015$~kg/m$^2$), because the departures $D$
would smooth are already tiny. $D$ moves neither coordinate of
Fig.~\ref{fig:regime}.

\section{Implications for V2}
\label{sec:implications}

\begin{center}
\begin{tabular}{llll}
\toprule
dial & sets & crossover & independent leverage \\
\midrule
$W_c$ & $\Wq$ level & $W_c^* = \Wstar - \tau_c E_0 = 43.91$ & none beyond $\Wq$ \\
$\tau_c$ & $\Wq$ offset $\tau_c E_0$ & $\tau_c^* = (\Wstar{-}W_c)/E_0$, iff $W_c{<}\Wstar$ & departure amplitude, window width \\
$\Delta_z$ & $\Wstar$ ($\propto$) & $\Delta_z^* = 68.20$~K at default $\Wq$ & dry circulation strength \\
$D$ & (neither) & none & eddy MSE flux magnitude \\
$v_d$ & (neither) & none & overturning $\Rightarrow$ departures, mean flux \\
\bottomrule
\end{tabular}
\end{center}

\begin{enumerate}
\item \textbf{Regime placement is now a closed-form choice.} To start V2 in
a chosen GMS regime, pick $\Wq - \Wstar$ directly; the boundary-adjacent
states already exist as runs with restart files ($W_c$=44 and
$\Delta_z$=68.2).
\item \textbf{$W_c$ and $\tau_c$ are degenerate for the regime.} Only their
combination $\Wq = W_c + \tau_c E_0$ decides the sign of $\Hhat$; they are
not independent regime dials. $\tau_c$'s independent role is the departure
amplitude (and, in V2, the stiffness of the precipitation coupling).
\item \textbf{V1's energetic blandness at the defaults is placement, not
structure.} At the default parameters the system sits $\sim$6~kg/m$^2$ deep
on the moisture-mode side, with the mean flux dominating the eddy flux
$\sim$25-fold. At the boundary that inverts: the eddy flux, the object the
ERA5-calibrated $D$ is meant to represent, becomes the entire MSE
transport.
\item \textbf{Expect the picture to deform under V2 coupling.} Everything
here is one-way; V2's latent heating lets $P$ anomalies drive the
circulation. On the $\Hhat < 0$ side that closes a positive feedback loop
(moisture-mode amplification), so the whole-domain-flips-as-a-unit behavior
and the narrow mixed-sign window of Sec.~\ref{sec:boundary} need not
survive coupling [speculative]. The near-boundary states are the natural V2 starting points
precisely because there the mean-flow feedback is weakest. The V2 gate
decision (kinematic $\Hv(\py v)$, validated dry) is already made.
\item \textbf{Costs.} A gated moist steady-state run is $\sim$1855 days,
measured at 27--45 min on the numpy backend ($\sim$0.9--1.4~s/day across
contention states); the numba kernel does not yet cover moisture (deferred
until V1 physics froze, which it now has).
\end{enumerate}

\appendix

\section{Model equations}
\label{app:eqs}

Upper-layer fields $u(y,t)$, $v(y,t)$, $\theta(y,t)$ on an equatorial
$\beta$-plane channel $|y| \le L$, walls at $\pm L$; moist V1 adds
$W(y,t)$. $\Hv$ is the Heaviside function.

\textbf{Zonal momentum} (SS09 Eq.~2.1 with the thermodynamic gate):
\begin{equation}
\frac{\partial u}{\partial t}
= v\big(\beta y - \py u\big)
\;-\; \Hv(\theta_E - \theta)\,(\py v)\,u
\;-\; \varepsilon_u u
\;-\; \underbrace{v_d\,\Hv(u)\,\mathrm{sgn}(y)\,\py u}_{\text{EMFD}} .
\end{equation}
The second term is vertical exchange with the lower layer where convection
is active (ascending air carries zero zonal momentum); the gate argument is
$\theta_E - \theta$ where SS09 use $\Hv(\py v)$, a deviation that suppresses
a grid-scale artifact of the original collocated grid (the two arguments
coincide in steady state through the thermodynamic balance; the kinematic
gate is validated on the staggered grid and will land with V2). The EMFD
(SS09 Eq.~2.5) represents extratropical eddies depositing easterly momentum
inside the westerly Hadley cell; it acts as poleward advection of $u$ at
speed $v_d$ where $u > 0$.

\textbf{Meridional momentum} (SS13 corrigendum form; the factor 2 is
intentional):
\begin{equation}
2\frac{\partial v}{\partial t} + \beta y\,u
= -\frac{gH}{T_0}\,\py T + k_v\,\partial_y^2 v,
\qquad T = \theta\,(p_t/p_s)^{R/c_p} = \theta/1.6 .
\end{equation}

\textbf{Thermodynamic}:
\begin{equation}
\frac{\partial \theta}{\partial t}
+ \frac{\delta\,\Delta_z}{H}\,\py v
= \frac{\theta_E - \theta}{\tau}
\;\big(+\,\kappa_\theta \partial_y^2 \theta,\ \kappa_\theta = 0\ \text{by
default}\big),
\end{equation}
with the equinoctial $\sin^2$ radiative-convective profile
\begin{equation}
\theta_E(y) = \theta_{00} - \Delta_y \sin^2\!\Big(\frac{\pi y}{2 y_1}\Big)
\quad (|y| < y_1), \qquad
\theta_E = \theta_{00} - \Delta_y \quad (|y| \ge y_1).
\end{equation}

\textbf{Moisture (V1, passive)}:
\begin{equation}
\frac{\partial W}{\partial t}
+ \py\big[\underbrace{-(2a{-}1)\,v\,W}_{\text{mean transport}}
\;\underbrace{-\,D\,\py W}_{\text{eddy mixing}}\big]
= E_0 - P,
\qquad P = \frac{(W - W_c)^+}{\tau_c}.
\end{equation}
The transport velocity is $-(2a{-}1)v$: the lower branch moves opposite the
upper-layer $v$, and with fraction $a$ of the column water in the lower
layer the column-integrated flux reduces to this coefficient. Nothing feeds
back on $u$, $v$, $\theta$.

\textbf{Column energetics} (derivations in
\path|docs/moist_v1_mse_budget.pdf|): with column heat capacity
$C = c_p \Pi \Delta p / g$ and $L_v = C\,\Lambda_{\mathrm{conv}}$,
\begin{equation}
\Shat = C\,\frac{\delta\,\Delta_z}{H}, \qquad
\Hhat = \Shat - L_v(2a{-}1)W, \qquad
\Wstar = \frac{\Shat}{L_v(2a{-}1)},
\end{equation}
and the steady column-MSE budget carries the three fluxes used throughout
this report: mean DSE $\Shat v$, mean latent $-L_v(2a{-}1)Wv$ (their sum is
the net mean flux $v\Hhat$), and eddy $-L_v D\,\py W$.

\textbf{Numerics.} Arakawa-C staggered grid ($v$ on the $\mathrm{ny}{-}1$
interior faces; $u$, $\theta$, $W$ on centers); leapfrog with Asselin
filter 0.04; upwind meridional advection of $u$; MC-limited one-sided EMFD
stencil (the gate-on production formulation); $W$ transport in
finite-volume flux form with MUSCL-MC face reconstruction (conserves total
water to roundoff); the moist diffusion $D$ and sink $P$ stepped at the
lagged $n{-}1$ level (stability bounds $2D\Delta t/\Delta y^2 \le 0.5$,
$2\Delta t/\tau_c \le 1$). Steady-state stopping: KE and $\theta$-variance
convergence plus the slow-drift gate (jet latitude, $|v|_{\max}$,
equatorial disequilibrium each flat to 0.2\% over a drag-timescale window).

\section{Parameter values}
\label{app:params}

Dry defaults (the Zhang et al.\ 2025 lineage; $k_v$ is $100\times$ SS09's
printed value, following their code):
\begin{center}
\begin{tabular}{llll}
\toprule
symbol & value & & meaning \\
\midrule
$\beta$ & $2\times10^{-11}$ & m$^{-1}$s$^{-1}$ & planetary vorticity gradient \\
$\varepsilon_u$ & $10^{-8}$ & s$^{-1}$ & Rayleigh drag on $u$ \\
$v_d$ & 2.5 & m/s & EMFD coefficient \\
$k_v$ & $7.786\times10^{5}$ & m$^2$/s & eddy viscosity on $v$ \\
$\tau$ & 37 & d & Newtonian cooling timescale \\
$H$ & 16\,000 & m & tropopause height \\
$\delta$ & 4\,000 & m & upper-layer thickness \\
$\Delta_z$ & 60 & K & vertical stability (scanned 30--90) \\
$T_0$ & 300 & K & reference temperature \\
$\theta_{00}$ & 330 & K & background $\theta_E$ \\
$\Delta_y$ & 50 & K & equator--pole $\theta_E$ contrast \\
$y_1$ & $9.439\times10^{6}$ & m & forcing half-width \\
$L$ & $1.5751\times10^{7}$ & m & channel half-width \\
$g$ & 9.81 & m/s$^2$ & gravity \\
\bottomrule
\end{tabular}
\end{center}

Moist V1 (\textbf{placeholders} pending the ERA5 calibration of $a$ and
$D$) and column constants:
\begin{center}
\begin{tabular}{llll}
\toprule
symbol & value & & meaning \\
\midrule
$a$ & 0.85 & & column-water-vapor fraction in the lower layer \\
$D$ & $10^{6}$ & m$^2$/s & eddy moisture diffusivity (scanned 0--2$\times10^6$) \\
$W_c$ & 50 & kg/m$^2$ & precipitation threshold (scanned 35--55) \\
$\tau_c$ & 14\,400 & s & convective timescale (scanned 1\,800--172\,800) \\
$E_0$ & $4.6\times10^{-5}$ & kg\,m$^{-2}$s$^{-1}$ & evaporation ($=3.97$ mm/d) \\
$C$ & $5.2\times10^{6}$ & J\,m$^{-2}$K$^{-1}$ & column heat capacity \\
$L_v$ & $2.5\times10^{6}$ & J/kg & latent heat ($= C\Lambda_{\mathrm{conv}}$) \\
\bottomrule
\end{tabular}
\end{center}

Numerics: $\mathrm{ny} = 801$ ($\Delta y = 39.4$~km), $\Delta t = 30$~s,
Asselin 0.04, numpy backend (numba does not yet cover moisture). Derived at
the defaults:
\begin{equation*}
\Shat = 7.8\times10^{7}~\mathrm{J/m^2},\quad
\Wstar = 44.57,\quad
\Wq = 50.66~\mathrm{kg/m^2},\quad
W_c^* = 43.91,\quad
\Delta_z^* = 68.20~\mathrm{K}.
\end{equation*}

\section{Run inventory and regeneration}
\label{app:runs}

All eighteen runs use
\begin{center}
\path|run-sw-model --ny 801 --dt 30 --enable-moisture --dw 1e6|\\
\path|--stop-at-steady-state --slow-drift-gate|
\end{center}
plus the dial flag(s)
(\path|--w-crit|, \path|--tau-c|, \path|--delta-z|), each in its own output
directory with its final restart file.

\begin{center}
\small
\begin{tabular}{lllr}
\toprule
scan & runs (directory) & dial values & gate day \\
\midrule
$D$ ladder & \path|moist_v1_validation/D{0,1,2}| & $D = \{0,1,2\}\times10^6$ & 1855 \\
$W_c$ scan & \path|moist_v1_wc_scan/wc{35,40,44,55}| & $W_c$ & 1855 \\
$\tau_c$ ladder & \path|moist_v1_tauc_dz/tc{1800,7200,57600,172800}| & $\tau_c$ & 1855 \\
$W_c$=40 pair & \path|moist_v1_tauc_dz/wc40tc{1800,172800}| & $W_c$, $\tau_c$ & 1855 \\
$\Delta_z$ ladder & \path|moist_v1_tauc_dz/dz{30,45,68p2,75,90}| & $\Delta_z$ & 1772--1896 \\
\bottomrule
\end{tabular}
\end{center}
(Directories relative to \path|model_output/|; every $\Delta_z$=60,
$\tau_c$=14\,400, $W_c$=50 slot reuses \path|moist_v1_validation/D1|. The
$\tau_c$ and $W_c$ runs share the D1 dry solution bit-for-bit, so gate day
1855 is exact; the $\Delta_z$ gate days are 1772, 1823, 1876, 1882, 1896
for 30, 45, 68.2, 75, 90~K.)

Analysis scripts (repo \path|scripts/|; each regenerates its figures from
the run directories): \path|moist_v1_analysis.py| ($D$ ladder + budget
figures), \path|moist_v1_wc_scan_analysis.py|,
\path|moist_v1_tauc_dz_analysis.py| (both ladders here, the scorecards,
Figs.~\ref{fig:tauc-ladder}--\ref{fig:dz-summary}), and
\path|moist_v1_regime_figure.py| (Fig.~\ref{fig:regime}). This report:
\path|latexmk -pdf -output-directory=docs docs/moist_v1_sensitivity_report.tex|
from the repo root.

\end{document}
